	\begin{center}
		\section{进阶数学}
	\end{center}

	\begin{multicols*}{2} % 开始双栏，数字 2 表示分两栏

		\subsection{FFT多项式乘法（非递归）}

		采用非递归版本实现FFT（蝴蝶变换），能显著减小常数并优化内存使用。
		这里通过 \texttt{multiply} 演示了如何运用 \texttt{fft} 对两个多项式（系数数组表示）进行乘法操作。

		\begin{minted}{cpp}
// 理论讲解 https://www.bilibili.com/video/BV1za411F76U/
// 参考（抄的） https://cp-algorithms.com/algebra/fft.html
// 多项式乘法 https://www.luogu.com.cn/record/229525769
using CD = complex<double>;

void fft(vector<CD> &a, bool invert) {
	int n = a.size();
	for (int i = 1, j = 0; i < n; i++) {
		int bit = n >> 1;
		for (; j & bit; bit >>= 1) j ^= bit;
		j ^= bit;
		if (i < j) swap(a[i], a[j]);
	}
	for (int len = 2; len <= n; len <<= 1) {
		double ang = 2 * pi / len * (invert ? -1 : 1);
		CD wlen(cos(ang), sin(ang));
		for (int i = 0; i < n; i += len) {
			CD w(1);
			for (int j = 0; j < len / 2; j++) {
				CD u = a[i + j], v = a[i + j + len / 2] * w;
				a[i + j] = u + v;
				a[i + j + len / 2] = u - v;
				w *= wlen;
			}
		}
	}
	if (invert) {
		for (CD &x : a) x /= n;
	}
}

vector<ll> multiply(vector<ll> &a, vector<ll> &b) {
	vector<CD> fa(a.begin(), a.end()), fb(b.begin(), b.end());
	int n = 1;
	while (n < (int)(a.size() + b.size())) {
		n <<= 1;
	}
	fa.resize(n); fb.resize(n);
	fft(fa, false); fft(fb, false);
	for (int i = 0; i < n; i++) fa[i] *= fb[i];
	fft(fa, true);
	vector<ll> res(n);
	// 避免精度损失
	for (int i = 0; i < n; i++) res[i] = llround(fa[i].real());
	return res;
}
		\end{minted}

		\subsubsection{高精度乘法（基于FFT）}

		https://www.luogu.com.cn/record/229526218

		大整数乘法的核心思想：将大整数的每一位看作多项式的一个系数。例如 $1234$ 可以表示为多项式 $4 + 3x + 2x^2 + x^3$（低位在前）。两个大整数相乘等价于两个多项式相乘，乘完后对结果数组统一处理进位。

		\begin{minted}{cpp}
inline void __() {
	string A, B;
	cin >> A >> B;
	vector<ll> a, b;
	for (int i = (int)A.size() - 1; i >= 0; --i) {
		a.push_back(A[i] - '0');
	}
	for (int i = (int)B.size() - 1; i >= 0; --i) {
		b.push_back(B[i] - '0');
	}
	
	vector<ll> ans = multiply(a, b);
	ll carry = 0;
	for (int i = 0; i < (int)ans.size(); ++i) {
		ans[i] += carry;
		carry = ans[i] / 10;
		ans[i] %= 10;
	}
	while (!ans.empty() && ans.back() == 0) {
		ans.pop_back();
	}
	if (ans.empty()) {
		cout << 0;
		return;
	}
	for (int i = (int)ans.size() - 1; i >= 0; --i) {
		cout << ans[i];
	}
}
		\end{minted}

		\subsubsection{FFT 的使用方法与常见应用}

		\texttt{multiply(a, b)} 的本质是：给定多项式 $A(x) = \sum a_i x^i$ 和 $B(x) = \sum b_i x^i$，返回乘积多项式 $C(x) = A(x) \cdot B(x)$ 的系数数组，其中 $c_k = \sum_{i+j=k} a_i \cdot b_j$。因此，\textbf{只要问题能转化为这种卷积形式，就能用 FFT 加速到 $O(n \log n)$}。

		\textbf{1. 加法卷积计数}

		最典型的场景：有两组数，问它们各取一个数相加能得到每个和的方案数。

		例如掷两颗骰子，求每种点数之和的方案数。将骰子 A 的各面出现次数存入 \texttt{a[]}（\texttt{a[k]} 表示点数 $k$ 出现的次数），骰子 B 同理存入 \texttt{b[]}，则 \texttt{multiply(a, b)} 的结果 \texttt{c[k]} 就是和为 $k$ 的方案数。

		\begin{minted}{cpp}
// 例：两个数组各取一个数，统计每种和的方案数
// a[i] = 数 i 在第一组中出现的次数
// b[j] = 数 j 在第二组中出现的次数
// c[k] = 和为 k 的方案数 = sum(a[i]*b[k-i])
vector<ll> c = multiply(a, b);
		\end{minted}

		\textbf{2. 字符串匹配（含通配符）}

		给定文本串 $T$（长度 $n$）和模式串 $P$（长度 $m$），对于 $T$ 的每个起始位置 $t$，定义失配度：
		$$D(t) = \sum_{j=0}^{m-1} (T_{t+j} - P_j)^2 \cdot w_j$$
		其中 $w_j$ 是通配符权重（通配符处 $w_j = 0$）。$D(t) = 0$ 表示位置 $t$ 完全匹配。

		将 $P$ 翻转为 $P'_j = P_{m-1-j}$，展开后各项都可化为卷积形式，分别用 FFT 计算后合并。

		\begin{minted}{cpp}
// 无通配符的简单情形：检查 T 和 P 是否在每个位置匹配
// 构造 a[i] = T[i], b[j] = P[m-1-j]（翻转 P）
// 则 c[t+m-1] = sum(T[t+j] * P[m-1-j]) 即为位置 t 的内积
// 更一般地，含通配符时需要构造多组卷积分别计算
vector<ll> a(n), b(m);
for (int i = 0; i < n; i++) a[i] = T[i] - 'a' + 1;
for (int j = 0; j < m; j++) b[j] = P[m - 1 - j] - 'a' + 1;
vector<ll> c = multiply(a, b);
// c[t + m - 1] 即为位置 t 的匹配内积
		\end{minted}

		\textbf{3. 高精度乘法}

		如上一节所述，将大整数每一位作为多项式系数，调用 \texttt{multiply} 后统一处理进位。

		\textbf{4. 多项式求值与插值的中间步骤}

		分治 FFT 可以加速多项式多点求值、多点插值等操作，核心都是递归地调用 \texttt{multiply} 合并子问题的结果多项式。

	\end{multicols*} % 从这里起改用单栏：多项式与线性规划的代码行都较长，双栏会频繁折行

		\subsection{NTT 卷积（ac-library 精简移植）}

		\href{https://www.luogu.com.cn/problem/P3803}{洛谷 P3803 多项式乘法} \\

		本节到「FWT 与子集卷积」为止的\textbf{全部代码}，都在 \texttt{template-check/src/poly\_check.cpp} 里
		与朴素实现随机对拍过（装了 ac-library 时还会与官方 NTT 逐点比对）；改动其中任何一段后请重跑它。

		上一节的复数 FFT 只能算整数卷积，一旦要在模 $998244353$ 意义下做多项式运算就得换成
		NTT。本节直接移植 \href{https://github.com/atcoder/ac-library}{ac-library} 的
		\texttt{convolution.hpp}，它是目前公开实现里常数最小的一档，理由有三：

		\begin{itemize}
			\item \textbf{radix-4 蝶形}：一次推进两层，乘法次数比常见的 radix-2 少约三成；
			\item \textbf{不需要位逆序置换数组}：正变换故意把结果留在位逆序上，逆变换再吃回去，
			      于是那个 \texttt{rev} 数组和它的预处理全都省掉了；
			\item \textbf{小规模自动降级}：短的一侧不超过 $60$ 时直接朴素乘，
			      避开两次变换的固定开销——多项式求逆、\texttt{exp} 这类倍增算法
			      会反复在很短的长度上调用卷积，这个降级的收益比想象中大。
		\end{itemize}

		\textbf{代价是 \texttt{butterfly} 的输出不是自然序}。只做点乘的场景
		（卷积、转置卷积、多点求值）完全不受影响；但如果哪天想在频域上「按下标做事」
		（比如手动取某一位的单位根），必须先想清楚下标已经被位逆序打乱了。

		\textbf{模数与原根}：NTT 要求模数形如 $c \cdot 2^k + 1$，且 $2^k$ 不小于变换长度。
		常用的几个如下（\texttt{rk} 即 $k$，决定了最长能做多长）：

		\begin{center}
		\begin{tabular}{rrrl}
			\hline
			模数 & 原根 & \texttt{rk} & 备注 \\
			\hline
			$998244353$   & $3$  & $23$ & 最常用，$= 119 \cdot 2^{23} + 1$ \\
			$1004535809   $ & $3$  & $21$ & 备用，与上者互质 \\
			$167772161$   & $3$  & $25$ & 三模数之一 \\
			$469762049$   & $3$  & $26$ & 三模数之一 \\
			$754974721$   & $\mathbf{11}$ & $24$ & 三模数之一，\textbf{原根不是 3} \\
			$2281701377$  & $3$  & $27$ & 超 \texttt{int} 范围，乘法要用 \texttt{\_\_int128} \\
			\hline
		\end{tabular}
		\end{center}

		最后一行的坑很值得记：$754974721$ 的原根是 $11$，抄板子时顺手写成 $3$，
		单模数的题照样能过（因为根本没用到它），一上三模数就会 WA 到怀疑人生。
		所以下面的模板把原根做成了\textbf{编译期自动求解}的默认模板参数，
		换模数时不用查表，也就没有写错的机会。

		\begin{minted}{cpp}
// ===== NTT 卷积：ac-library atcoder/convolution.hpp 精简移植 =====
// 参考（抄的） https://github.com/atcoder/ac-library/blob/master/atcoder/convolution.hpp
// 模板题 https://www.luogu.com.cn/problem/P3803
// 编译期快速幂 + 最小原根：换模数时不用查表，也就不会写错原根
constexpr ull powConst(ull b, ull e, ull m) {
    ull r = 1;
    for (b %= m; e; e >>= 1, b = b * b % m)
        if (e & 1) r = r * b % m;
    return r;
}
constexpr uint32_t primitiveRoot(uint32_t m) {
    uint32_t divs[20] = {2}, x = (m - 1) / 2;
    int cnt = 1;
    while (x % 2 == 0) x /= 2;
    for (uint32_t i = 3; (ull)i * i <= x; i += 2)
        if (x % i == 0) {
            divs[cnt++] = i;
            while (x % i == 0) x /= i;
        }
    if (x > 1) divs[cnt++] = x;
    for (uint32_t g = 2;; g++) {
        bool ok = true;
        for (int i = 0; i < cnt && ok; i++)
            ok = powConst(g, (m - 1) / divs[i], m) != 1;
        if (ok) return g;
    }
}

// MOD 必须形如 c * 2^k + 1（k 要够大），G 默认取最小原根
template <uint32_t MOD, uint32_t G = primitiveRoot(MOD)>
struct NTT {
    using mint = modint<uint32_t, MOD>;
    static constexpr int rk = __builtin_ctz(MOD - 1);   // 最长支持 2^rk
    struct Info {
        mint root[rk + 1], iroot[rk + 1];               // root[i]^(2^i) == 1
        mint rate2[rk - 1], irate2[rk - 1];
        mint rate3[rk - 2], irate3[rk - 2];
        Info() {
            root[rk] = mint::powmod(mint(G), (MOD - 1) >> rk);
            iroot[rk] = mint::inv(root[rk]);
            for (int i = rk - 1; i >= 0; i--) {
                root[i] = root[i + 1] * root[i + 1];
                iroot[i] = iroot[i + 1] * iroot[i + 1];
            }
            mint prod = 1, iprod = 1;
            for (int i = 0; i + 2 <= rk; i++) {
                rate2[i] = root[i + 2] * prod;
                irate2[i] = iroot[i + 2] * iprod;
                prod *= iroot[i + 2], iprod *= root[i + 2];
            }
            prod = 1, iprod = 1;
            for (int i = 0; i + 3 <= rk; i++) {
                rate3[i] = root[i + 3] * prod;
                irate3[i] = iroot[i + 3] * iprod;
                prod *= iroot[i + 3], iprod *= root[i + 3];
            }
        }
    };
    static const Info &info() { static Info x; return x; }

    // 正变换：一次推两层（radix-4），出来是位逆序，因此不需要 rev 置换数组
    static void butterfly(vector<mint> &a) {
        const Info &w = info();
        int n = a.size(), h = __builtin_ctz((unsigned)n), len = 0;
        while (len < h) {
            if (h - len == 1) {                       // 只剩一层，走 radix-2
                int p = 1 << (h - len - 1);
                mint rot = 1;
                for (int s = 0; s < (1 << len); s++) {
                    int off = s << (h - len);
                    for (int i = 0; i < p; i++) {
                        mint l = a[i + off], r = a[i + off + p] * rot;
                        a[i + off] = l + r, a[i + off + p] = l - r;
                    }
                    if (s + 1 != (1 << len))
                        rot *= w.rate2[__builtin_ctz(~(unsigned)s)];
                }
                len++;
            } else {                                  // radix-4
                int p = 1 << (h - len - 2);
                mint rot = 1, imag = w.root[2];
                for (int s = 0; s < (1 << len); s++) {
                    mint rot2 = rot * rot, rot3 = rot2 * rot;
                    int off = s << (h - len);
                    for (int i = 0; i < p; i++) {
                        ull m2 = (ull)MOD * MOD;
                        ull a0 = a[i + off].val();
                        ull a1 = (ull)a[i + off + p].val() * rot.val();
                        ull a2 = (ull)a[i + off + 2 * p].val() * rot2.val();
                        ull a3 = (ull)a[i + off + 3 * p].val() * rot3.val();
                        ull t = (ull)mint(a1 + m2 - a3).val() * imag.val();
                        ull na2 = m2 - a2;
                        a[i + off] = a0 + a2 + a1 + a3;
                        a[i + off + p] = a0 + a2 + (2 * m2 - (a1 + a3));
                        a[i + off + 2 * p] = a0 + na2 + t;
                        a[i + off + 3 * p] = a0 + na2 + (m2 - t);
                    }
                    if (s + 1 != (1 << len))
                        rot *= w.rate3[__builtin_ctz(~(unsigned)s)];
                }
                len += 2;
            }
        }
    }
    // 逆变换：吃位逆序输入，还原自然序（未除以 n，交给 conv 统一处理）
    static void butterflyInv(vector<mint> &a) {
        const Info &w = info();
        int n = a.size(), h = __builtin_ctz((unsigned)n), len = h;
        while (len) {
            if (len == 1) {
                int p = 1 << (h - len);
                mint irot = 1;
                for (int s = 0; s < (1 << (len - 1)); s++) {
                    int off = s << (h - len + 1);
                    for (int i = 0; i < p; i++) {
                        mint l = a[i + off], r = a[i + off + p];
                        a[i + off] = l + r;
                        a[i + off + p] =
                            (ull)((uint32_t)(l.val() - r.val()) + MOD) * irot.val();
                    }
                    if (s + 1 != (1 << (len - 1)))
                        irot *= w.irate2[__builtin_ctz(~(unsigned)s)];
                }
                len--;
            } else {
                int p = 1 << (h - len);
                mint irot = 1, iimag = w.iroot[2];
                for (int s = 0; s < (1 << (len - 2)); s++) {
                    mint irot2 = irot * irot, irot3 = irot2 * irot;
                    int off = s << (h - len + 2);
                    for (int i = 0; i < p; i++) {
                        ull a0 = a[i + off].val(), a1 = a[i + off + p].val();
                        ull a2 = a[i + off + 2 * p].val();
                        ull a3 = a[i + off + 3 * p].val();
                        ull t = mint((MOD + a2 - a3) * iimag.val()).val();
                        a[i + off] = a0 + a1 + a2 + a3;
                        a[i + off + p] = (a0 + (MOD - a1) + t) * irot.val();
                        a[i + off + 2 * p] =
                            (a0 + a1 + (MOD - a2) + (MOD - a3)) * irot2.val();
                        a[i + off + 3 * p] =
                            (a0 + (MOD - a1) + (MOD - t)) * irot3.val();
                    }
                    if (s + 1 != (1 << (len - 2)))
                        irot *= w.irate3[__builtin_ctz(~(unsigned)s)];
                }
                len -= 2;
            }
        }
    }
    static vector<mint> naive(const vector<mint> &a, const vector<mint> &b) {
        vector<mint> c(a.size() + b.size() - 1);
        for (size_t i = 0; i < a.size(); i++)
            for (size_t j = 0; j < b.size(); j++) c[i + j] += a[i] * b[j];
        return c;
    }
    // 短的一侧不超过 60 时直接朴素乘，省掉两次变换的常数
    static vector<mint> conv(vector<mint> a, vector<mint> b) {
        int n = a.size(), m = b.size();
        if (!n || !m) return {};
        if (min(n, m) <= 60) return naive(a, b);
        int z = 1;
        while (z < n + m - 1) z <<= 1;
        assert((MOD - 1) % z == 0);   // 长度超过 2^rk：换模数或走三模数版本
        a.resize(z), b.resize(z);
        butterfly(a), butterfly(b);
        for (int i = 0; i < z; i++) a[i] *= b[i];
        butterflyInv(a);
        a.resize(n + m - 1);
        mint iz = mint::inv(mint(z));
        for (auto &x : a) x *= iz;
        return a;
    }
};

constexpr uint32_t P = 998244353;   // 原根 3，rk = 23（最长做到 2^23）
using Z = modint<uint32_t, P>;
vector<Z> convolution(const vector<Z> &a, const vector<Z> &b) {
    return NTT<P>::conv(a, b);
}
		\end{minted}

		\textbf{用法}：\texttt{convolution(a, b)} 返回 $c_k = \sum_{i+j=k} a_i b_j$，
		长度是 \texttt{a.size() + b.size() - 1}。依赖基础数学一章的自动取模类 \texttt{modint}。

		\begin{minted}{cpp}
vector<Z> a{1, 2, 3}, b{4, 5};       // (1+2x+3x^2)(4+5x)
vector<Z> c = convolution(a, b);     // c = {4, 13, 22, 15}
// 换模数只要改这两行，原根会自己算出来：
// constexpr uint32_t P = 1004535809;
// using Z = modint<uint32_t, P>;
		\end{minted}

		\subsection{多项式全家桶（Poly）}

		把系数数组包成一个继承自 \texttt{vector<Z>} 的结构体，所有操作都建立在上面那个
		\texttt{convolution} 之上。约定：下标即次数，\texttt{f[0]} 是常数项；
		所有带参数 \texttt{m} 的运算都只保证\textbf{前 $m$ 项正确}
		（幂级数运算本来就是在 $\bmod\ x^m$ 意义下做的）。

		\textbf{基础部分}：取前 $k$ 项、平移、加减、数乘、求导、积分。
		\texttt{get(i)} 越界返回 $0$，这样加减法不用先对齐长度。

		\begin{minted}{cpp}
// ===== 多项式全家桶：系数下标从低到高，f = a[0] + a[1] x + a[2] x^2 + ... =====
struct Poly : vector<Z> {
    using vector<Z>::vector;
    Poly(const vector<Z> &a) : vector<Z>(a) {}
    Z get(int i) const { return i < 0 || i >= (int)size() ? Z(0) : (*this)[i]; }
    int deg() const {              // 最高非零次数，零多项式返回 -1
        int i = (int)size() - 1;
        while (i >= 0 && (*this)[i] == Z(0)) i--;
        return i;
    }
    Poly &norm() { resize(deg() + 1); return *this; }
    Poly trunc(int k) const {      // 取前 k 项，不足补零
        Poly f = *this;
        f.resize(k);
        return f;
    }
    Poly shift(int k) const {      // 乘上 x^k，k 为负则右移（丢掉低位）
        if (k >= 0) {
            Poly f = *this;
            f.insert(f.begin(), k, Z(0));
            return f;
        }
        if ((int)size() <= -k) return Poly();
        return Poly(begin() - k, end());
    }
    friend Poly operator+(const Poly &a, const Poly &b) {
        Poly c(max(a.size(), b.size()));
        for (int i = 0; i < (int)c.size(); i++) c[i] = a.get(i) + b.get(i);
        return c;
    }
    friend Poly operator-(const Poly &a, const Poly &b) {
        Poly c(max(a.size(), b.size()));
        for (int i = 0; i < (int)c.size(); i++) c[i] = a.get(i) - b.get(i);
        return c;
    }
    friend Poly operator-(const Poly &a) {
        Poly c = a;
        for (auto &x : c) x = -x;
        return c;
    }
    friend Poly operator*(const Poly &a, Z k) {
        Poly c = a;
        for (auto &x : c) x *= k;
        return c;
    }
    friend Poly operator*(Z k, const Poly &a) { return a * k; }
    friend Poly operator*(const Poly &a, const Poly &b) {
        return Poly(convolution(a, b));
    }
    Poly &operator+=(const Poly &b) { return *this = *this + b; }
    Poly &operator-=(const Poly &b) { return *this = *this - b; }
    Poly &operator*=(const Poly &b) { return *this = *this * b; }
    Poly &operator*=(Z k) { return *this = *this * k; }
		\end{minted}

		\textbf{牛顿迭代五件套}。它们的骨架是同一个：设已经求出 $\bmod\ x^k$ 的解 $x_0$，
		要把精度翻倍到 $\bmod\ x^{2k}$，就对方程 $G(x) = 0$ 做一步牛顿法
		$$x \equiv x_0 - \frac{G(x_0)}{G'(x_0)} \pmod{x^{2k}}.$$
		代入不同的 $G$ 就得到不同的运算：求逆取 $G(x) = 1/x - f$，
		开方取 $G(x) = x^2 - f$。因为每一步的规模翻倍，总复杂度是
		$T(n) = T(n/2) + O(n\log n) = O(n\log n)$——\textbf{和一次卷积同阶}，
		这也是为什么多项式科技能在 $10^5$ 级别的题上跑得动。

		\begin{minted}{cpp}
    Poly deriv() const {           // 求导
        if (empty()) return Poly();
        Poly f(size() - 1);
        for (int i = 1; i < (int)size(); i++) f[i - 1] = (*this)[i] * Z(i);
        return f;
    }
    Poly integr() const {          // 积分，常数项取 0
        Poly f(size() + 1);
        for (int i = 0; i < (int)size(); i++) f[i + 1] = (*this)[i] / Z(i + 1);
        return f;
    }
    // 乘法逆，要求 a[0] != 0    模板题 https://www.luogu.com.cn/problem/P4238
    Poly inv(int m) const {
        Poly x{Z::inv((*this)[0])};
        for (int k = 1; k < m; k <<= 1)
            x = (x * (Poly{Z(2)} - trunc(k << 1) * x)).trunc(k << 1);
        return x.trunc(m);
    }
    // ln，要求 a[0] == 1        模板题 https://www.luogu.com.cn/problem/P4725
    Poly log(int m) const { return (deriv() * inv(m)).integr().trunc(m); }
    // exp，要求 a[0] == 0       模板题 https://www.luogu.com.cn/problem/P4726
    Poly exp(int m) const {
        Poly x{Z(1)};
        for (int k = 1; k < m; k <<= 1)
            x = (x * (Poly{Z(1)} - x.log(k << 1) + trunc(k << 1))).trunc(k << 1);
        return x.trunc(m);
    }
    // 开方，这里只处理 a[0] == 1；a[0] 是别的二次剩余时先解出 sqrt(a[0]) 再整体缩放
    // 模板题 https://www.luogu.com.cn/problem/P5205
    Poly sqrt(int m) const {
        Poly x{Z(1)};
        Z inv2 = Z::inv(Z(2));
        for (int k = 1; k < m; k <<= 1)
            x = (x + (trunc(k << 1) * x.inv(k << 1)).trunc(k << 1)) * inv2;
        return x.trunc(m);
    }
    // k 次幂，允许最低非零位不在 0 处、也允许 a[i] != 1
    // 模板题 https://www.luogu.com.cn/problem/P5245
    Poly pow(ll k, int m) const {
        if (k == 0) { Poly r(m); if (m) r[0] = 1; return r; }
        int i = 0;
        while (i < (int)size() && (*this)[i] == Z(0)) i++;
        if (i == (int)size() || (__int128)i * k >= m) return Poly(m);
        Z v = (*this)[i];
        Poly f = shift(-i) * Z::inv(v);
        int rest = m - i * (int)k;
        return (f.log(rest) * Z(k % (ll)P)).exp(rest).shift(i * (int)k)
               * Z::powmod(v, k);
    }
		\end{minted}

		\textbf{使用前提别记错}：\texttt{inv} 要 $f_0 \neq 0$；\texttt{log} 要 $f_0 = 1$
		（否则 $\ln f$ 的常数项不存在）；\texttt{exp} 要 $f_0 = 0$
		（否则 $e^{f}$ 的常数项是个超越数）；\texttt{sqrt} 这份只处理 $f_0 = 1$，
		$f_0$ 是别的二次剩余时先用 Cipolla 解出 $\sqrt{f_0}$ 再整体缩放。
		\texttt{pow} 已经处理好了「最低非零项不在 $0$ 次」的情况：先把 $x^i$ 提出来，
		剩下的部分归一化成 $f_0 = 1$ 再走 $\exp(k \ln f)$。

		\textbf{转置卷积与多点求值}。\texttt{mulT} 取 $f \cdot \operatorname{rev}(b)$ 的高位，
		它算的是「相关」而不是卷积；多点求值就是拿它在分治树上把
		$\prod (1 - x_i t)$ 一层层往下推，$O(n \log^2 n)$。

		\begin{minted}{cpp}
    // 转置卷积：取 (*this) * reverse(b) 的高位，多点求值的核心零件
    Poly mulT(const Poly &b) const {
        if (b.empty()) return Poly();
        Poly c = b;
        reverse(c.begin(), c.end());
        return (*this * c).shift(-((int)b.size() - 1));
    }
    // 多点求值 O(n log^2 n)   模板题 https://www.luogu.com.cn/problem/P5050
    vector<Z> eval(vector<Z> x) const {
        if (empty()) return vector<Z>(x.size(), Z(0));
        int n = max((int)x.size(), (int)size());
        vector<Poly> q(n << 2);
        vector<Z> ans(x.size());
        x.resize(n);
        auto build = [&](auto &self, int p, int l, int r) -> void {
            if (r - l == 1) { q[p] = Poly{Z(1), -x[l]}; return; }
            int m = (l + r) >> 1;
            self(self, p << 1, l, m), self(self, p << 1 | 1, m, r);
            q[p] = q[p << 1] * q[p << 1 | 1];
        };
        build(build, 1, 0, n);
        auto work = [&](auto &self, int p, int l, int r, const Poly &f) -> void {
            if (r - l == 1) { if (l < (int)ans.size()) ans[l] = f.get(0); return; }
            int m = (l + r) >> 1;
            self(self, p << 1, l, m, f.mulT(q[p << 1 | 1]).trunc(m - l));
            self(self, p << 1 | 1, m, r, f.mulT(q[p << 1]).trunc(r - m));
        };
        work(work, 1, 0, n, mulT(q[1].inv(n)));
        return ans;
    }
};
		\end{minted}

		\textbf{带余除法}。设 $\deg f = n$、$\deg g = m$，令 $f^R$ 表示系数翻转，
		由 $f = qg + r$ 两边翻转可得
		$$f^R \equiv q^R g^R \pmod{x^{n-m+1}},$$
		于是 $q^R = f^R \cdot (g^R)^{-1} \bmod x^{n-m+1}$，再翻回去、回代求 $r$。

		\begin{minted}{cpp}
// 带余除法 f = q * g + r，deg(r) < deg(g)
// 模板题 https://www.luogu.com.cn/problem/P4512
pair<Poly, Poly> divmod(Poly f, Poly g) {
    int n = f.deg(), m = g.deg();
    f.norm(), g.norm();
    if (n < m) return {Poly(), f};
    Poly fr = f, gr = g;
    reverse(fr.begin(), fr.end()), reverse(gr.begin(), gr.end());
    Poly q = (fr.trunc(n - m + 1) * gr.inv(n - m + 1)).trunc(n - m + 1);
    reverse(q.begin(), q.end());
    Poly r = (f - q * g).trunc(max(m, 1));
    return {q.norm(), r.norm()};
}
		\end{minted}

		\textbf{快速插值}是多点求值的对偶：给 $n$ 个点值还原出 $n-1$ 次多项式。
		记 $M(x) = \prod (x - x_i)$，则
		$$f(x) = \sum_i \frac{y_i}{M'(x_i)} \prod_{j \neq i} (x - x_j),$$
		分母那一堆 $M'(x_i)$ 用一次多点求值批量算出来，分子在分治树上自底向上合并。

		\begin{minted}{cpp}
// 快速插值：给 n 个点值还原 n-1 次多项式 O(n log^2 n)
// 模板题 https://www.luogu.com.cn/problem/P5158
Poly interpolate(const vector<Z> &x, const vector<Z> &y) {
    int n = x.size();
    vector<Poly> m(n << 2);
    auto build = [&](auto &self, int p, int l, int r) -> void {
        if (r - l == 1) { m[p] = Poly{-x[l], Z(1)}; return; }
        int mid = (l + r) >> 1;
        self(self, p << 1, l, mid), self(self, p << 1 | 1, mid, r);
        m[p] = m[p << 1] * m[p << 1 | 1];
    };
    build(build, 1, 0, n);
    vector<Z> d = m[1].deriv().eval(x);      // d[i] = M'(x_i)
    auto solve = [&](auto &self, int p, int l, int r) -> Poly {
        if (r - l == 1) return Poly{y[l] / d[l]};
        int mid = (l + r) >> 1;
        return self(self, p << 1, l, mid) * m[p << 1 | 1]
             + self(self, p << 1 | 1, mid, r) * m[p << 1];
    };
    return solve(solve, 1, 0, n);
}
		\end{minted}

		\begin{center}
		\begin{tabular}{lll}
			\hline
			操作 & 复杂度 & 模板题 \\
			\hline
			\texttt{operator*} & $O(n\log n)$ & P3803 \\
			\texttt{inv} & $O(n\log n)$ & P4238 \\
			\texttt{log} / \texttt{exp} & $O(n\log n)$ & P4725 / P4726 \\
			\texttt{sqrt} & $O(n\log n)$ & P5205 \\
			\texttt{pow} & $O(n\log n)$ & P5245 \\
			\texttt{divmod} & $O(n\log n)$ & P4512 \\
			\texttt{eval} & $O(n\log^2 n)$ & P5050 \\
			\texttt{interpolate} & $O(n\log^2 n)$ & P5158 \\
			\hline
		\end{tabular}
		\end{center}

		\subsection{任意模数卷积（三模数 NTT + Garner）}

		\href{https://www.luogu.com.cn/problem/P4245}{洛谷 P4245 任意模数多项式乘法} \\

		题目给的模数是 $10^9+7$ 这类\textbf{不支持 NTT} 的数时，常见有两条路：

		\begin{itemize}
			\item \textbf{拆系数 FFT（MTT）}：把每个系数拆成 $a = a_1 \cdot 2^{15} + a_0$，
			      做三到四次复数 FFT 再合并。代码短，但全程靠 \texttt{double} 的 53 位尾数硬扛，
			      系数和长度一起上去就会掉精度，而且掉了不报错。
			\item \textbf{三模数 NTT + Garner}：在三个 NTT 友好的模数下各算一遍，
			      再用 CRT 把真值还原出来，最后对题目模数取模。
			      \textbf{全程整数运算，没有精度问题}，代价是常数约为单次 NTT 的三倍。
		\end{itemize}

		这里取后者。三个模数的乘积约 $5.9 \times 10^{25}$，足够容纳
		「长度 $10^6$、系数 $10^9$」量级卷积的真值（上界约 $10^{24}$）。

		\begin{minted}{cpp}
// ===== 任意模数 / 大整数卷积：三模数 NTT + Garner 合并 =====
// 模板题 https://www.luogu.com.cn/problem/P4245
// 三个模数乘积 ~5.9e25，足够容纳任意 int 系数、长度 1e6 级别的卷积真值
// 注意原根各不相同（167772161 与 469762049 是 3，754974721 是 11），
// 这里靠 primitiveRoot 编译期算出来，不必手写
constexpr uint32_t M1 = 167772161, M2 = 469762049, M3 = 754974721;

template <uint32_t MOD>
vector<ll> convRaw(const vector<ll> &a, const vector<ll> &b) {
    using mint = modint<uint32_t, MOD>;
    vector<mint> A(a.size()), B(b.size());
    for (size_t i = 0; i < a.size(); i++) A[i] = a[i];
    for (size_t i = 0; i < b.size(); i++) B[i] = b[i];
    auto C = NTT<MOD>::conv(A, B);
    vector<ll> c(C.size());
    for (size_t i = 0; i < C.size(); i++) c[i] = C[i].val();
    return c;
}

// p > 0：结果对 p 取模，p 可以是 1e9+7 这种非 NTT 模数
// p <= 0：不取模，要求真实值 < 2^63（否则请自行换成 __int128 返回）
vector<ll> convolutionAnyMod(const vector<ll> &a, const vector<ll> &b, ll p) {
    if (a.empty() || b.empty()) return {};
    auto c1 = convRaw<M1>(a, b), c2 = convRaw<M2>(a, b), c3 = convRaw<M3>(a, b);
    ll i12 = modint<uint32_t, M2>::inv(modint<uint32_t, M2>(M1)).val();
    ll i3 = modint<uint32_t, M3>::inv(
                modint<uint32_t, M3>((ll)((ull)M1 * M2 % M3))).val();
    vector<ll> c(c1.size());
    for (size_t i = 0; i < c1.size(); i++) {
        ll k1 = (c2[i] - c1[i]) % (ll)M2 * i12 % (ll)M2;
        if (k1 < 0) k1 += M2;
        ll x12 = c1[i] + (ll)M1 * k1;                  // < M1*M2 ~ 7.9e16
        ll k2 = (c3[i] - x12) % (ll)M3 * i3 % (ll)M3;
        if (k2 < 0) k2 += M3;
        __int128 x = (__int128)x12 + (__int128)M1 * M2 * k2;
        c[i] = p > 0 ? (ll)(x % p) : (ll)x;
    }
    return c;
}
		\end{minted}

		\textbf{两个用途}：\texttt{p > 0} 时是任意模数卷积；\texttt{p <= 0} 时不取模，
		直接给出真实的整数结果（要求真值 $< 2^{63}$），可以当作\textbf{高精度乘法}
		或「大整数卷积」用——比上一节的复数 FFT 版本精度可靠得多。

		\subsection{分治 FFT}

		\href{https://www.luogu.com.cn/problem/P4721}{洛谷 P4721 分治 FFT} \\

		要解的是这类自指递推：$f_0$ 已知，且
		$$f_i = \sum_{j=1}^{i} f_{i-j}\, g_j .$$
		麻烦在于 $f$ 的每一项都依赖前面所有项，不能一次卷积算完。CDQ 分治的思路是
		「先把左半边算对，再把左半边对右半边的贡献一次卷积推过去」。

		\begin{minted}{cpp}
// ===== 分治 FFT：f[0] 已知，f[i] = sum_{j=1..i} f[i-j] * g[j] =====
// 模板题 https://www.luogu.com.cn/problem/P4721
// 也可以直接写成 f = f0 / (1 - g)，即 Poly{1} - g 求逆，常数更小、代码更短
vector<Z> divideConquerFFT(const vector<Z> &g, int n, Z f0 = 1) {
    vector<Z> f(n);
    if (!n) return f;
    f[0] = f0;
    auto cdq = [&](auto &self, int l, int r) -> void {
        if (r - l == 1) return;
        int m = (l + r) >> 1;
        self(self, l, m);
        vector<Z> a(f.begin() + l, f.begin() + m);
        vector<Z> b(g.begin(), g.begin() + min((int)g.size(), r - l));
        vector<Z> c = convolution(a, b);
        for (int i = m; i < r; i++)
            if (i - l < (int)c.size()) f[i] += c[i - l];
        self(self, m, r);
    };
    cdq(cdq, 0, n);
    return f;
}
		\end{minted}

		\textbf{但更该记住的是它的等价形式}。把递推写成生成函数：
		$F = F G + f_0$，于是
		$$F = \frac{f_0}{1 - G}.$$
		直接一次多项式求逆就完了，代码只有一行，常数还比分治小：

		\begin{minted}{cpp}
// 与 divideConquerFFT(g, n) 等价，但只需一次求逆
Poly f = (Poly{Z(1)} - Poly(g)).inv(n) * f0;
		\end{minted}

		那还留着分治版做什么？因为\textbf{当 $g$ 本身依赖 $f$ 时（两个数列互相牵扯），
		生成函数写不出封闭形式，分治仍然是唯一可行的写法}。
		「先看看能不能写成求逆」应该是遇到这类递推的第一反应，写不成再退回分治。

		\subsection{线性递推：Berlekamp-Massey 与 Bostan-Mori}

		\href{https://www.luogu.com.cn/problem/P5487}{洛谷 P5487 BM} ·
		\href{https://www.luogu.com.cn/problem/P4723}{洛谷 P4723 常系数线性递推} \\

		这两个板子是一对，合起来能做一件很赖皮的事：
		\textbf{打出前几十项，直接问第 $10^{18}$ 项是多少}。

		\textbf{Berlekamp-Massey} 从一段数列里找出最短的线性递推式。经验口径：
		若怀疑答案满足 $k$ 阶递推，喂给它至少 $2k$ 项（保险起见 $2k + 10$），
		返回的系数 $c$ 满足 $s_i = \sum_j c_j\, s_{i-1-j}$。

		\begin{minted}{cpp}
// ===== Berlekamp-Massey：求最短线性递推式 O(n^2) =====
// 模板题 https://www.luogu.com.cn/problem/P5487
// 返回 c，满足 s[i] = sum_{j} c[j] * s[i-1-j]
vector<Z> berlekampMassey(const vector<Z> &s) {
    vector<Z> ls, cur;
    int lf = 0;
    Z ld = 0;
    for (int i = 0; i < (int)s.size(); i++) {
        Z t = 0;
        for (int j = 0; j < (int)cur.size(); j++) t += cur[j] * s[i - 1 - j];
        if (t == s[i]) continue;
        if (cur.empty()) {                    // 还没有任何递推式
            cur.resize(i + 1);
            lf = i, ld = t - s[i];
            continue;
        }
        Z k = (t - s[i]) / ld;
        vector<Z> c(i - lf - 1);
        c.push_back(k);
        for (int j = 0; j < (int)ls.size(); j++) c.push_back(-ls[j] * k);
        if (c.size() < cur.size()) c.resize(cur.size());
        for (int j = 0; j < (int)cur.size(); j++) c[j] += cur[j];
        if (i - (int)cur.size() > lf - (int)ls.size()) ls = cur, lf = i, ld = t - s[i];
        cur = c;
    }
    return cur;
}
		\end{minted}

		\textbf{Bostan-Mori} 求有理分式的第 $n$ 项系数 $[x^n]\dfrac{P(x)}{Q(x)}$。
		核心恒等式是分母有理化：
		$$\frac{P(x)}{Q(x)} = \frac{P(x)Q(-x)}{Q(x)Q(-x)},$$
		右边的分母 $Q(x)Q(-x)$ 是 $x^2$ 的多项式，于是按 $n$ 的奇偶性抽取分子的
		奇/偶次项，问题规模从 $n$ 降到 $n/2$，共 $O(\log n)$ 轮。

		\begin{minted}{cpp}
// ===== Bostan-Mori：求 [x^n] P(x)/Q(x)，O(k log k log n) =====
Z bostanMori(Poly p, Poly q, ll n) {
    while (n) {
        Poly qm = q;                          // Q(-x)
        for (int i = 1; i < (int)qm.size(); i += 2) qm[i] = -qm[i];
        Poly u = p * qm, v = q * qm;          // V(x^2) = Q(x)Q(-x)，只剩偶次项
        Poly np(u.size() / 2 + 1), nq(v.size() / 2 + 1);
        for (int i = (int)(n & 1); i < (int)u.size(); i += 2) np[i >> 1] = u[i];
        for (int i = 0; i < (int)v.size(); i += 2) nq[i >> 1] = v[i];
        p = np, q = nq, n >>= 1;
    }
    return p.get(0) / q[0];
}

// 常系数齐次线性递推：给出前 k 项 a[0..k-1] 与系数 c（BM 的输出可直接喂进来），
// 求 a[n]。模板题 https://www.luogu.com.cn/problem/P4723
Z linearRecurrence(const vector<Z> &a, const vector<Z> &c, ll n) {
    int k = c.size();
    if (n < (ll)a.size()) return a[n];
    Poly q(k + 1);
    q[0] = 1;
    for (int i = 0; i < k; i++) q[i + 1] = -c[i];
    Poly p = (Poly(a) * q).trunc(k);
    return bostanMori(p, q, n);
}
		\end{minted}

		\textbf{典型用法}：矩阵快速幂只能做到 $O(k^3\log n)$，$k$ 一大就没救；
		这套是 $O(k\log k\log n)$，$k = 10^4$、$n = 10^{18}$ 也能秒。

		\begin{minted}{cpp}
vector<Z> s = {1, 1, 2, 3, 5, 8, 13, 21};   // 打表出来的前几项
vector<Z> c = berlekampMassey(s);           // c = {1, 1}，即 f[i]=f[i-1]+f[i-2]
Z ans = linearRecurrence(s, c, 1000000000000000000LL);
		\end{minted}

		\subsection{拉格朗日插值}

		\href{https://www.luogu.com.cn/problem/P4781}{洛谷 P4781 拉格朗日插值} \\

		已知 $n$ 个点值求 $n-1$ 次多项式在某点的取值：
		$$f(k) = \sum_{i} y_i \prod_{j \neq i} \frac{k - x_j}{x_i - x_j}.$$

		\begin{minted}{cpp}
// ===== 拉格朗日插值 =====
// 任意 n 个点值求 f(k)，O(n^2)  模板题 https://www.luogu.com.cn/problem/P4781
Z lagrange(const vector<Z> &x, const vector<Z> &y, Z k) {
    int n = x.size();
    Z ans = 0;
    for (int i = 0; i < n; i++) {
        Z num = y[i], den = 1;
        for (int j = 0; j < n; j++)
            if (j != i) num *= k - x[j], den *= x[i] - x[j];
        ans += num / den;
    }
    return ans;
}
		\end{minted}

		\textbf{横坐标取连续整数 $0,1,\dots,n-1$ 时能降到 $O(n)$}：分母是
		$i!\,(n-1-i)!$ 带个符号，分子 $\prod_{j\neq i}(k-j)$ 用前后缀积。
		这是竞赛里最常用的形态：

		\begin{minted}{cpp}
// 横坐标是 0,1,...,n-1 时可以做到 O(n)：前后缀积 + 阶乘逆元
// 典型用途：自然数幂和 sum_{i=1..m} i^t 是 t+1 次多项式，取 t+2 个点值即可
Z lagrangeConsecutive(const vector<Z> &y, ll k) {
    int n = y.size();
    if (k < n) return y[k];
    vector<Z> pre(n + 1), suf(n + 1), fac(n), ifac(n);
    pre[0] = suf[n] = 1;
    for (int i = 0; i < n; i++) pre[i + 1] = pre[i] * (Z(k) - Z(i));
    for (int i = n - 1; i >= 0; i--) suf[i] = suf[i + 1] * (Z(k) - Z(i));
    fac[0] = 1;
    for (int i = 1; i < n; i++) fac[i] = fac[i - 1] * Z(i);
    ifac[n - 1] = Z::inv(fac[n - 1]);
    for (int i = n - 1; i >= 1; i--) ifac[i - 1] = ifac[i] * Z(i);
    Z ans = 0;
    for (int i = 0; i < n; i++) {
        Z t = y[i] * pre[i] * suf[i + 1] * ifac[i] * ifac[n - 1 - i];
        ans += ((n - 1 - i) & 1) ? -t : t;
    }
    return ans;
}
		\end{minted}

		\textbf{最典型的应用是自然数幂和}。$\sum_{i=1}^{m} i^t$ 关于 $m$ 是一个
		$t+1$ 次多项式（可由伯努利数或差分法证），所以\textbf{只要暴力算出前 $t+2$ 个点值，
		插值就能得到任意 $m$ 处的答案}，$m$ 可以大到 $10^{18}$：

		\begin{minted}{cpp}
// 求 sum_{i=1..m} i^t （t 很小，m 很大）
vector<Z> y(t + 2);
for (int i = 1; i <= t + 1; i++) y[i] = y[i - 1] + Z::powmod(Z(i), t);
Z ans = lagrangeConsecutive(y, m);
		\end{minted}

		\subsection{FWT 与子集卷积}

		\href{https://www.luogu.com.cn/problem/P4717}{洛谷 P4717 FWT} ·
		\href{https://www.luogu.com.cn/problem/P6097}{洛谷 P6097 子集卷积} \\

		FFT 处理的是\textbf{加法卷积} $c_k = \sum_{i+j=k} a_i b_j$；
		把「$i+j$」换成位运算就是 FWT 的地盘：
		$$c_k = \sum_{i \,\mathrm{op}\, j = k} a_i b_j, \qquad
		  \mathrm{op} \in \{\,\mathrm{or},\ \mathrm{and},\ \mathrm{xor}\,\}.$$
		三者的变换矩阵不同，但骨架完全一样：按位分治，每层把
		$(a_j,\ a_{j+len})$ 换成两个线性组合。或卷积做的是子集和（高维前缀和），
		与卷积是超集和，异或卷积则是每一位上的 $\pm$ 组合。

		\begin{minted}{cpp}
// ===== FWT：位运算卷积 =====
// 模板题 https://www.luogu.com.cn/problem/P4717
// type: 0 = 或(OR)  1 = 与(AND)  2 = 异或(XOR)；数组长度必须是 2 的幂
void fwt(vector<Z> &a, int type, bool inv) {
    int n = a.size();
    for (int len = 1; len < n; len <<= 1)
        for (int i = 0; i < n; i += len << 1)
            for (int j = i; j < i + len; j++) {
                Z u = a[j], v = a[j + len];
                if (type == 0) a[j + len] = inv ? v - u : v + u;
                else if (type == 1) a[j] = inv ? u - v : u + v;
                else a[j] = u + v, a[j + len] = u - v;
            }
    if (type == 2 && inv) {
        Z iv = Z::inv(Z(n));
        for (auto &x : a) x *= iv;
    }
}

vector<Z> bitwiseConv(vector<Z> a, vector<Z> b, int type) {
    fwt(a, type, false), fwt(b, type, false);
    for (size_t i = 0; i < a.size(); i++) a[i] *= b[i];
    fwt(a, type, true);
    return a;
}
		\end{minted}

		\textbf{子集卷积}要求的是\textbf{不交并}：$c_S = \sum_{T \subseteq S} a_T\, b_{S\setminus T}$。
		直接用或卷积会把 $T$ 与 $S\setminus T$ 有交的情况一起算进去。
		经典的修补办法是\textbf{占位多项式}：给每个集合额外记一维 $\mathrm{popcount}$，
		只有当 $|T| + |S \setminus T| = |S|$ 时才是真正的不交并，
		于是在「或卷积」的基础上让这一维做加法卷积即可，复杂度 $O(n^2 2^n)$。

		\begin{minted}{cpp}
// ===== 子集卷积：c[S] = sum_{T subset S} a[T] * b[S\T]，O(n^2 2^n) =====
// 模板题 https://www.luogu.com.cn/problem/P6097
// 做法：给或卷积补一维「popcount」，占位后维度相加才是真正的不交并
vector<Z> subsetConv(const vector<Z> &a, const vector<Z> &b, int n) {
    int N = 1 << n;
    vector<vector<Z>> fa(n + 1, vector<Z>(N)), fb(n + 1, vector<Z>(N));
    vector<vector<Z>> fc(n + 1, vector<Z>(N));
    for (int s = 0; s < N; s++) {
        fa[__builtin_popcount(s)][s] = a[s];
        fb[__builtin_popcount(s)][s] = b[s];
    }
    for (int i = 0; i <= n; i++) fwt(fa[i], 0, false), fwt(fb[i], 0, false);
    for (int i = 0; i <= n; i++)
        for (int j = 0; i + j <= n; j++)
            for (int s = 0; s < N; s++) fc[i + j][s] += fa[i][s] * fb[j][s];
    for (int i = 0; i <= n; i++) fwt(fc[i], 0, true);
    vector<Z> c(N);
    for (int s = 0; s < N; s++) c[s] = fc[__builtin_popcount(s)][s];
    return c;
}
		\end{minted}

		\subsection{拉格朗日反演}

		\href{https://www.luogu.com.cn/problem/P5825}{洛谷 P5825 排列计数}（典型应用） \\

		生成函数题里经常出现这种局面：\textbf{$G$ 只满足一个方程、写不出封闭形式}，
		但它的\textbf{复合逆} $F$（即 $F(G(x)) = x$）反而形式简单。
		拉格朗日反演说，这时候不必真的把 $G$ 解出来：
		\[
			[x^n]\,G(x) = \frac1n\,[x^{n-1}]\left(\frac{x}{F(x)}\right)^{n},
		\]
		更常用的是带一个外层函数的扩展形式
		\[
			[x^n]\,H(G(x)) = \frac1n\,[x^{n-1}]\,H'(x)\left(\frac{x}{F(x)}\right)^{n}.
		\]
		成立条件是 $F_0 = 0$ 且 $F_1 \neq 0$（这正是复合逆存在的条件）。

		\begin{minted}{cpp}
// ===== 拉格朗日反演：由 F 直接取出它复合逆 G 的某一项系数 =====
// F(G(x)) = x，要求 F[0] = 0 且 F[1] != 0
// 模板题 https://www.luogu.com.cn/problem/P5825 （常用于 GF 只写得出方程的场合）
// [x^n] G = (1/n) [x^{n-1}] (x / F(x))^n
Z lagrangeInversion(const Poly &f, int n) {
    Poly t = f.shift(-1).inv(n);            // t = x / F(x)
    return t.pow(n, n).get(n - 1) / Z(n);
}
// 扩展形式：[x^n] H(G(x)) = (1/n) [x^{n-1}] H'(x) (x / F(x))^n
Z lagrangeInversion(const Poly &f, const Poly &h, int n) {
    Poly t = f.shift(-1).inv(n);
    return (h.deriv() * t.pow(n, n)).get(n - 1) / Z(n);
}
		\end{minted}

		\textbf{怎么用}：设未知的计数序列是 $G$，从组合意义列出它满足的方程
		（常见形如 $G = x\,\Phi(G)$），把它改写成 $F(G) = x$ 的样子读出 $F$，
		再套上面的式子。例如卡特兰数：由 $G = x + G^2$ 得 $F(x) = x - x^2$，于是
		\[
			[x^n]G = \frac1n [x^{n-1}](1-x)^{-n} = \frac1n\binom{2n-2}{n-1},
		\]
		正是 $C_{n-1}$。这条路子的价值在于：\textbf{哪怕 $\Phi$ 复杂到 $G$ 根本解不出来，
		只要能写出 $F$，单点系数依然是 $O(n\log n)$ 可算的}。

		\subsection{生成函数常用结论}

		代码只解决「算得快」，真正卡人的是\textbf{把计数问题翻译成幂级数}这一步。
		这一节是给赛场上现查的对照表。

		\textbf{两种生成函数}。普通型（OGF）与指数型（EGF）：
		$$F(x) = \sum_{n \ge 0} a_n x^n, \qquad
		  \hat F(x) = \sum_{n \ge 0} a_n \frac{x^n}{n!} .$$
		选哪个取决于\textbf{对象之间有没有编号}：组合出来的结构\textbf{无标号}
		（如整数拆分、括号串）用 OGF，因为两个 OGF 相乘正好是「左边取一部分、右边取一部分」；
		\textbf{有标号}（如排列、带编号的图）用 EGF，因为
		$$\hat F \hat G = \sum_n \left(\sum_{k} \binom{n}{k} f_k\, g_{n-k}\right)\frac{x^n}{n!}$$
		里的 $\binom{n}{k}$ 恰好完成了「把 $n$ 个编号分给两边」这件事。

		\textbf{OGF 常见封闭形式}：

		\begin{itemize}
			\item $\displaystyle \sum_{n\ge 0} x^n = \frac{1}{1-x}$，
			      更一般地 $\displaystyle \sum_{n\ge 0} x^{kn} = \frac{1}{1-x^k}$
			      （每种物品可取任意多个）；
			\item $\displaystyle 1 + x + \cdots + x^{m-1} = \frac{1-x^m}{1-x}$
			      （每种物品最多取 $m-1$ 个）；
			\item $\displaystyle \sum_{n\ge 0}(n+1)x^n = \frac{1}{(1-x)^2}$，
			      一般地 $\displaystyle \sum_{n \ge 0}\binom{n+m}{m}x^n = \frac{1}{(1-x)^{m+1}}$；
			\item $\displaystyle \sum_{n\ge 0}\binom{m}{n}x^n = (1+x)^m$（二项式定理）；
			\item \textbf{广义二项式定理}：对任意实数 $\alpha$，
			      $\displaystyle (1+x)^{\alpha} = \sum_{n \ge 0}\binom{\alpha}{n}x^n$，
			      其中 $\binom{\alpha}{n} = \dfrac{\alpha(\alpha-1)\cdots(\alpha-n+1)}{n!}$；
			\item $\displaystyle \sum_{n \ge 1}\frac{x^n}{n} = -\ln(1-x)$；
			\item \textbf{斐波那契}：$\displaystyle \sum_{n\ge0}F_n x^n
			      = \frac{(F_1-F_0)x + F_0}{1-x-x^2}$
			      （由 $G = xG + x^2G + (F_1-F_0)x + F_0$ 解出）；
			\item \textbf{卡特兰数}：$\displaystyle \sum_{n\ge0}C_n x^n
			      = \frac{1-\sqrt{1-4x}}{2x}$（由 $C = xC^2 + 1$ 解出）；
			\item \textbf{前缀和算子}：若 $s_n = \sum_{i\le n} a_i$，则
			      $\displaystyle S(x) = \frac{F(x)}{1-x}$；乘 $1-x$ 就是差分；
			\item \textbf{五边形数定理}：$\displaystyle \prod_{i\ge1}(1-x^i)
			      = \sum_{k\in\mathbb{Z}}(-1)^k x^{k(3k-1)/2}$。
			      分拆数 $p(n)$ 的 OGF 是它的倒数 $\prod (1-x^i)^{-1}$，
			      于是 $p(n)$ 可以用这个只有 $O(\sqrt n)$ 项的稀疏级数
			      在 $O(n\sqrt n)$ 内递推，不需要多项式求逆。
		\end{itemize}

		\textbf{EGF 与 $\exp$ 的组合意义}。这是 EGF 最有威力的地方：
		$$\hat F = \exp \hat P \iff
		  \text{「}F\text{ 是若干个无序的 }P\text{ 拼起来」}.$$
		直观理解是 $\exp \hat P = \sum_k \hat P^k / k!$：取 $k$ 个 $P$ 结构，
		除以 $k!$ 消掉它们之间的顺序。反过来取 $\ln$ 就是「从任意结构里剥出连通结构」。
		由此可得一串对偶：

		\begin{itemize}
			\item $\displaystyle \sum_{n\ge0}\frac{x^n}{n!} = e^x$（每个元素只有一种选择）；
			\item 长度为 $n$ 的\textbf{轮换}有 $(n-1)!$ 种，
			      $\displaystyle \hat P = \sum_{n\ge1}\frac{x^n}{n} = -\ln(1-x)$；
			      而排列是若干轮换的无序拼接，故排列数的 EGF
			      $\displaystyle \exp\hat P = \frac{1}{1-x}$，即 $n!$，自洽；
			\item 限定轮换长度 $\ge 2$（即无不动点）得 \textbf{错排}：
			      $\hat P = -\ln(1-x) - x$，错排数的 EGF 为
			      $\displaystyle \exp\hat P = \frac{e^{-x}}{1-x}$；
			\item $n$ 个点的\textbf{有标号无向图}有 $2^{\binom n2}$ 个，其 EGF 记为 $\hat G$；
			      \textbf{连通图}的 EGF 就是 $\ln \hat G$——
			      这是「连通图计数」标准做法的全部内容；
			\item $n$ 个点的\textbf{有标号树}有 $n^{n-2}$ 个（Cayley 公式），
			      森林的 EGF 即 $\exp$ 掉它。
		\end{itemize}

		\textbf{单位根反演}。要从卷积结果里只取下标为 $k$ 倍数的项：
		$$[k \mid n] = \frac{1}{k}\sum_{j=0}^{k-1}\omega_k^{jn},
		  \qquad \omega_k = e^{2\pi i/k}\ \ (\text{模意义下取 } g^{(p-1)/k}).$$
		于是 $\displaystyle \sum_{k \mid n} a_n = \frac1k \sum_{j=0}^{k-1} F(\omega_k^j)$。
		$k$ 很小（$2$、$3$、$4$）时这招最好使，常见于「计数模 $k$ 余 $r$ 的方案」。

		\textbf{解题套路}：把「每个物品能选多少个」写成一个多项式，
		所有物品的多项式\textbf{相乘}就是总方案的生成函数；
		若干个同类物品用 $\exp$（EGF，无序）或 $\frac{1}{1-x^k}$（OGF）合并；
		最后要么化成封闭形式直接展开，要么用本章的多项式全家桶把它算出来。

		\subsection{线性规划（两阶段单纯形 Simplex）}

		求解标准型
		$$
		\max\ c^{\mathsf T} z \quad \text{s.t.} \quad A z \le b,\ z \ge 0 .
		$$
		三样输入分别是：$A$ 是\textbf{约束系数}（$A_{ij}$ = 第 $i$ 条约束里变量 $j$ 的系数），$b$ 是每条约束的\textbf{上限}，$c$ 是\textbf{要最大化的那个式子的系数}。代码里就直接叫 \texttt{coef} / \texttt{limit} / \texttt{profit}，解向量叫 \texttt{x}。

		可行域是若干半空间交出的凸多面体，最优值必在某个\textbf{顶点}处取到。单纯形法就是从一个顶点出发，每次沿一条棱走到相邻的更优顶点（一次 \texttt{pivot} = 一个非基变量入基、一个基变量出基），直到目标行里没有改进方向为止。

		\begin{center}
		\begin{tikzpicture}[scale=1.5, every node/.style={font=\small}]
			\coordinate (O) at (0,0);
			\coordinate (A) at (3,0);
			\coordinate (B) at (3,1);
			\coordinate (C) at (1.5,2.5);
			\coordinate (D) at (0,3);
			% 约束直线（延伸到可行域外用浅色虚线）
			\draw[gray!55, dashed] (0,4) -- (4,0);
			\draw[gray!55, dashed] (0,3) -- (4,1.667);
			\draw[gray!55, dashed] (3,-0.25) -- (3,3.3);
			\fill[blue!9] (O) -- (A) -- (B) -- (C) -- (D) -- cycle;
			\draw[blue!60, thick] (O) -- (A) -- (B) -- (C) -- (D) -- cycle;
			\draw[->] (-0.25,0) -- (4.5,0) node[below] {$x$};
			\draw[->] (0,-0.25) -- (0,3.8) node[left] {$y$};
			% 约束名
			\node[font=\scriptsize, gray!70!black] at (4.05,0.35) {$x+y\le 4$};
			\node[font=\scriptsize, gray!70!black] at (4.05,1.95) {$x+3y\le 9$};
			\node[font=\scriptsize, gray!70!black] at (3.0,3.6) {$x\le 3$};
			% 单纯形游走路径
			\draw[->, red!80!black, very thick] (O) -- (A);
			\draw[->, red!80!black, very thick] (A) -- (B);
			\foreach \p in {O,A,B,C,D} \fill (\p) circle (1.7pt);
			\node[font=\scriptsize, below left=1pt and 1pt] at (O) {$z=0$};
			\node[font=\scriptsize, below=5pt] at (A) {$z=9$};
			\node[font=\scriptsize, right=4pt, red!70!black] at (B) {$z=11$};
			\node[font=\scriptsize, above right=1pt and 2pt] at (C) {$z=9.5$};
			\node[font=\scriptsize, above left=1pt and 2pt] at (D) {$z=6$};
		\end{tikzpicture}
		\end{center}

		图中目标为 $\max z = 3x + 2y$，路径 $(0,0) \to (3,0) \to (3,1)$ 即两次 pivot 后到达最优顶点。

		\textbf{实现要点}：

		\begin{itemize}
			\item \textbf{字典表连续存储}：$(rows+2) \times (cols+2)$ 一维数组行主序，比 \texttt{vector<vector<double>>} 少一层间接寻址，内层循环无分支可向量化。
			\item \textbf{两阶段}：$b$ 有负分量时原点不可行，先引入人工变量 $x_0$ 跑第一阶段 $\min x_0$；若最优 $x_0 > 0$ 则原问题\textbf{不可行}。$b \ge 0$ 时直接跳过第一阶段。
			\item \textbf{Devex 定价}：入基列挑 $d_j^2 / w_j$ 最大者（$d_j$ 为目标行系数、$w_j$ 为参考权重），比朴素 Dantzig 规则（挑最负的 $d_j$）迭代次数少很多。
			\item \textbf{防循环}：退化时目标值长期不推进，连续 \texttt{stallLimit} 次后切 Bland 最小下标规则（保证有限步终止），推进后自动切回 Devex。
			\item \textbf{出基行}：最小比值 $b_i / a_{is}$；比值持平时选主元绝对值更大的一行，数值更稳。
		\end{itemize}

		\textbf{返回值}：最优值；无界返回 $+\infty$，不可行返回 $-\infty$；\texttt{solution} 非空时回填一组最优解。

		\textbf{建模转换}：等式约束拆成 $\le$ 与 $\ge$ 两条；$\ge$ 约束两边取负变 $\le$；$\min$ 问题把 \texttt{profit} 整体取负、答案再取负；变量可正可负时拆成 $z = z^{+} - z^{-}$ 两个非负变量。

		\begin{minted}{cpp}
// 两个变量 x、y，都 >= 0；要最大化 3x + 2y
// 三条约束：x + y <= 4，x + 3y <= 9，x <= 3
vector<vector<double>> coef = {{1, 1},      // 第 1 条约束里 x、y 的系数
                               {1, 3},      // 第 2 条
                               {1, 0}};     // 第 3 条（y 不出现，系数写 0）
vector<double> limit = {4, 9, 3};           // 每条约束不等号右边的上限
vector<double> profit = {3, 2};             // 目标 3x + 2y 里 x、y 的系数
SimplexLP lp;
lp.init(coef, limit, profit);
vector<double> x;                           // 取到最优时每个变量的值
double best = lp.solve(&x);					// best = 11, x = {3, 1}
		\end{minted}

		\begin{minted}{cpp}
// 变量 x[0..cols) 全都 >= 0；每条约束形如 coef[i][0]*x[0] + ... <= limit[i]；
// 目标是让 profit[0]*x[0] + profit[1]*x[1] + ... 尽量大。
struct SimplexLP {
    // 核心就是一张表 tab，(rows + 2) 行 x (cols + 2) 列，拍平成一维存（比二维 vector 少一层寻址）：
    //   每行一条约束，每列一个变量，最后一列放这条约束还剩多少余量
    //   多出的倒数第 2 行 = 目标（每个变量再加一单位能赚多少），最后一行只在"先找可行点"时用
    //   多出的倒数第 2 列 = 辅助变量 x0，同样只在"先找可行点"时用
    int rows = 0;               // 约束条数
    int cols = 0;               // 变量个数
    int stride = 0;             // 一行占多少个 double（= cols + 2）
    vector<double> tab;
    vector<int> basis;          // basis[i] = 第 i 行当前解出来的那个变量的编号
    vector<int> nonbasis;       // nonbasis[j] = 第 j 列对应的变量编号，这些变量当前取 0
    double eps = 1e-9;
    long long pivots = 0;       // 换基次数，调试看规模用
    vector<double> weight;      // Devex 参考权重，挑变量时用

    double* row(int i) { return tab.data() + static_cast<size_t>(i) * stride; }
    const double* row(int i) const { return tab.data() + static_cast<size_t>(i) * stride; }

    // coef[i][j]：第 i 条约束里变量 j 的系数     limit[i]：第 i 条约束的上限
    // profit[j] ：变量 j 每多一单位，目标增加多少
    void init(const vector<vector<double>>& coef,
              const vector<double>& limit,
              const vector<double>& profit) {
        rows = static_cast<int>(limit.size());
        cols = static_cast<int>(profit.size());
        stride = cols + 2;
        tab.assign(static_cast<size_t>(rows + 2) * stride, 0.0);
        basis.resize(rows);
        nonbasis.resize(cols + 1);
        for (int i = 0; i < rows; ++i) {
            double* p = row(i);
            for (int j = 0; j < cols; ++j) p[j] = coef[i][j];
            p[cols] = -1.0;          // 辅助变量 x0 那一列，只在找可行点时起作用
            p[cols + 1] = limit[i];
            basis[i] = cols + i;     // 一开始所有变量取 0，这一行剩的余量就是 limit[i]
        }
        double* obj = row(rows);
        // 目标行存的是 -profit：以后看到负数就说明"这个变量还能加，目标还能变大"
        for (int j = 0; j < cols; ++j) { obj[j] = -profit[j]; nonbasis[j] = j; }
        nonbasis[cols] = -1;
        row(rows + 1)[cols] = 1.0;   // 找可行点那一阶段的目标：把 x0 压到 0
        weight.assign(cols + 1, 1.0);
        pivots = 0;
    }

    // 让第 s 列的变量从 0 开始往大加，加到第 r 条约束把它卡住为止；这两个变量交换身份
    void pivot(int r, int s) {
        ++pivots;
        double* pr = row(r);
        const double inv = 1.0 / pr[s];
        {   // Devex：顺手更新参考权重，供下次挑变量用（O(cols)）
            const double as = pr[s], ws = weight[s];
            for (int j = 0; j <= cols; ++j) {
                if (j == s) continue;
                const double ratio = pr[j] / as;
                const double cand = ratio * ratio * ws;
                if (cand > weight[j]) weight[j] = cand;
            }
            const double nw = ws / (as * as);
            weight[s] = nw > 1.0 ? nw : 1.0;
        }
        const int total = rows + 2, width = cols + 2;
        for (int i = 0; i < total; ++i) {   // 用第 r 行把其余每行的第 s 列消成 0
            if (i == r) continue;
            double* pi = row(i);
            const double f = pi[s] * inv;
            if (f == 0.0) continue;
            for (int j = 0; j < width; ++j) pi[j] -= pr[j] * f;   // 内层无分支，编译器可向量化
            pi[s] = -f;                                           // 第 s 列要单独修正
        }
        for (int j = 0; j < width; ++j) pr[j] *= inv;
        pr[s] = inv;
        swap(basis[r], nonbasis[s]);
    }

    // phase = 1：先找一个满足所有约束的点（把 x0 压到 0）
    // phase = 2：在可行的范围内把目标做到最大
    bool run(int phase) {
        const int objRow = (phase == 1) ? rows + 1 : rows;
        // 目标值长时间不动（退化）就换 Bland 规则：它慢，但保证不会绕圈死循环
        long long stall = 0;
        const long long stallLimit = 4LL * (rows + cols) + 64;
        while (true) {
            const double* obj = row(objRow);
            int s = -1;                                  // 先挑"加谁"：哪个变量还能让目标变大
            if (stall < stallLimit) {                    // Devex：挑性价比最高的
                double best = 0.0;
                for (int j = 0; j <= cols; ++j) {
                    if (phase == 2 && nonbasis[j] == -1) continue;
                    const double d = obj[j];
                    if (d > -eps) continue;              // 不是负的 => 加它目标也不会变大
                    const double score = d * d / weight[j];
                    if (score > best) { best = score; s = j; }
                }
            } else {                                     // Bland：老实挑编号最小的
                for (int j = 0; j <= cols; ++j) {
                    if (phase == 2 && nonbasis[j] == -1) continue;
                    if (obj[j] < -eps && (s == -1 || nonbasis[j] < nonbasis[s])) s = j;
                }
            }
            if (s == -1) return true;                    // 没变量能让目标变大 => 这一阶段做完

            // 再看"加到哪停"：哪条约束最先被顶满，即余量 / 系数 最小的那一行
            int r = -1;
            double bestRatio = 0.0, bestPiv = 0.0;
            for (int i = 0; i < rows; ++i) {
                const double* pi = row(i);
                if (pi[s] < eps) continue;               // 系数 <= 0，这条约束拦不住它
                const double ratio = pi[cols + 1] / pi[s];
                if (r == -1 || ratio < bestRatio - 1e-12 ||
                    // 两行一样紧时选系数绝对值更大的那行，算起来数值更稳
                    (ratio < bestRatio + 1e-12 && pi[s] > bestPiv)) {
                    r = i; bestRatio = ratio; bestPiv = pi[s];
                }
            }
            if (r == -1) return false;                   // 没人拦得住它 => 目标能到无穷大

            const double before = row(objRow)[cols + 1];
            pivot(r, s);
            const double after = row(objRow)[cols + 1];
            stall = (after > before + 1e-12) ? 0 : stall + 1;
        }
    }

    // 返回目标能达到的最大值；无界返回 +inf，一个可行点都没有返回 -inf
    // 传了 solution 就顺便把最优时每个变量的取值填回去
    double solve(vector<double>* solution = nullptr) {
        int r = 0;
        for (int i = 1; i < rows; ++i)
            if (row(i)[cols + 1] < row(r)[cols + 1]) r = i;
        if (row(r)[cols + 1] < -eps) {                   // 有约束上限是负的 => 全取 0 都不合法，先找可行点
            pivot(r, cols);
            if (!run(1) || row(rows + 1)[cols + 1] < -eps)
                return -numeric_limits<double>::infinity();   // x0 压不到 0 => 无解
            for (int i = 0; i < rows; ++i) {
                if (basis[i] != -1) continue;            // 把还赖在解里的 x0 换出去
                int s = -1;
                for (int j = 0; j <= cols; ++j)
                    if (s == -1 || row(i)[j] < row(i)[s] ||
                        (row(i)[j] == row(i)[s] && nonbasis[j] < nonbasis[s])) s = j;
                pivot(i, s);
            }
        }
        if (!run(2)) return numeric_limits<double>::infinity();
        if (solution) {                                  // 没被解出来的变量取值就是 0
            solution->assign(cols, 0.0);
            for (int i = 0; i < rows; ++i)
                if (basis[i] < cols) (*solution)[basis[i]] = row(i)[cols + 1];
        }
        return row(rows)[cols + 1];
    }
};
		\end{minted}

